\begin{abstract}
This work introduces \textsc{QLoRA}, a finetuning technique designed to dramatically reduce memory demands, enabling the finetuning of models with up to 65B parameters on a single 48GB GPU, while retaining the full efficacy of standard 16-bit finetuning. \textsc{QLoRA} operates by propagating gradients through a 4-bit quantized, frozen pretrained language model into Low Rank Adapters~(LoRA). Our primary model suite, \textbf{Guanaco}, sets a new bar among open models on the Vicuna benchmark, attaining 99.3\% of ChatGPT’s performance with just 24 hours of finetuning on one GPU. The method brings several key innovations that enable substantial memory conservation without compromising accuracy: (a) the introduction of 4-bit NormalFloat (NF4), a novel data format optimized from an information-theoretic standpoint for normally distributed weights; (b) Double Quantization, which compresses the quantization parameters themselves to further minimize memory usage; and (c) Paged Optimizers, which efficiently handle temporary increases in memory usage. Using \textsc{QLoRA}, we have finetuned over 1,000 models and conduct a comprehensive assessment of instruction-following and chatbot capabilities across eight instruction datasets, various architectures (including LLaMA and T5), and large model sizes that traditional finetuning cannot practically handle (such as 33B and 65B parameters). Our experiments indicate that QLoRA-finetuned models, even at modest sizes, can surpass prior state-of-the-art results when trained on a curated high-quality dataset. We supplement our empirical results with an extensive evaluation using both human raters and GPT-4, showing that GPT-4-based evaluations offer a reliable and cost-effective substitute for human judgment. We also identify reliability issues with existing chatbot benchmarks in faithfully reflecting real chatbot performance. A targeted analysis highlights specific limitations of \textbf{Guanaco} relative to ChatGPT. All models, codebase, and CUDA kernels for 4-bit training are openly available.\footnote{\url{https://github.com/artidoro/qlora} and \url{https://github.com/TimDettmers/bitsandbytes}}
\end{abstract}

\section{Introduction}

The process of finetuning large language models (LLMs) has proven to be a powerful approach for enhancing model capabilities \citep{min2021metaicl, wei2021finetuned, ouyang2022training, wang2022super, wang2022self, liu2022few} as well as for instilling or eliminating specific model behaviors \citep{ouyang2022training,askell2021general,bai2022training}. Nevertheless, updating parameters in extremely large-scale models poses a considerable financial and computational burden; for example, standard 16-bit finetuning for LLaMA with 65 billion parameters~\citep{touvron2023llama} demands over 780 GB of GPU memory. Although recent advances in quantization significantly lower LLMs' memory usage during inference \citep{dettmers2022llm, dettmers2022case, frantar2022gptq, xiao2022smoothquant}, these solutions are generally incompatible with training, resulting in instability or failure \citep{wortsman2023stable}. 

In this work, we introduce the first method enabling finetuning of quantized 4-bit models without loss in model performance. Our proposed approach, \textsc{QLoRA}, employs a unique high-precision quantization scheme that reduces a pretrained model to 4 bits, then supplements it with a compact set of trainable Low-rank Adapter weights \citep{hu2021lora}, 

which are optimized by propagating gradients through the quantized parameters during training. 

With \textsc{QLoRA}, memory demands for finetuning a 65B-parameter model are lowered from more than 780 GB to under 48 GB of GPU memory, while maintaining both inference speed and accuracy when compared to the traditional 16-bit finetuning paradigm. This development dramatically enhances the practicality and availability of finetuning large public models, enabling single-GPU adaptation of the largest open models currently accessible. Utilizing \textsc{QLoRA}, we present the \textbf{Guanaco} series of models, of which our second largest achieves 97.8\% of ChatGPT’s score on the Vicuna~\citep{vicuna2023} benchmark, while completing finetuning in under 12 hours using one consumer GPU. Employing a professional-grade GPU over 24 hours allows our most capable model to reach 99.3\% of ChatGPT’s performance, effectively matching ChatGPT on the Vicuna benchmark. In deployment scenarios, the smallest variant in the \textbf{Guanaco} lineup (7B parameters) requires merely 5 GB of memory and outperforms the 26 GB Alpaca model by a margin exceeding 20 percentage points on the Vicuna evaluation (see Table~\ref{tbl:vicuna_eval}).

\textsc{QLoRA} incorporates several novel strategies aimed at minimizing memory requirements while maintaining model effectiveness: (1) \textbf{4-bit NormalFloat}, a quantization format optimized for data with normal distributions, which outperforms both 4-bit Integer and 4-bit Float types in practice. (2) \textbf{Double Quantization}, which entails quantizing the quantization parameters themselves, thereby reducing memory usage by about 0.37 bits per parameter—translating to around 3 GB saved for a 65B-parameter model. (3) \textbf{Paged Optimizers}, leveraging NVIDIA unified memory technology to eliminate memory spikes associated with gradient checkpointing when processing mini-batches containing long sequences. Together, these techniques are integrated into a refined LoRA framework with adapter layers present throughout the entire network, effectively mitigating the accuracy compromises noted in earlier approaches.

Due to the improved efficiency of \textsc{QLoRA}, it becomes feasible to conduct a comprehensive examination of instruction finetuning and chatbot performance at model scales previously unattainable because of prohibitive memory demands. To this end, we train in excess of 1,000 models, covering a broad array of instruction tuning datasets, model families, and parameter counts ranging from 80M to 65B. Alongside demonstrations that \textsc{QLoRA} matches the performance of standard 16-bit training (\S\ref{sec:comp_to_std}) and enables the training of a leading chatbot, \textbf{Guanaco}, (\S\ref{sec:pushing}), our analysis reveals notable patterns. Most notably, the quality of training data emerges as considerably more critical than the sheer quantity; for instance, the OASST1 dataset, containing only 9,000 samples, outperforms the 450,000-sample (subsampled) FLAN v2 dataset in chatbot metrics, even though both aim to foster generalization in instruction following. Additionally, we show that excelling on the Massive Multitask Language Understanding (MMLU) benchmark does not necessarily translate to high performance on the Vicuna chatbot benchmark (and vice versa), highlighting that the appropriateness of the dataset to the target task outweighs the size for that domain.

In addition, we conduct a thorough assessment of chatbot effectiveness, employing both evaluations from human raters and GPT-4. Our methodology incorporates a tournament-based benchmarking framework, wherein different chatbot models compete head-to-head by generating responses to identical prompts. The determination of which model produces the superior response in each pairwise contest is carried out by either human judges or GPT-4. Cumulative outcomes from these matches are then converted into Elo ratings~\citep{elo1967proposed, elo1978rating}, which serve as a metric to order the models by performance. Our findings reveal a substantial concordance between GPT-4 and human raters regarding the overall performance hierarchy, though notable divergences also occur. Thus, we caution that while model-driven assessment provides a cost-effective alternative to human judgment, it introduces additional uncertainties.

We complement our benchmarking results with an in-depth qualitative examination of the \textbf{Guanaco} family of models, documenting instances of both notable strengths and evident failures that were not apparent from quantitative evaluation alone.

To support future investigations, we have released the full set of model outputs along with both human and GPT-4 annotations. Our software—including code, CUDA kernels, and integration with the Hugging Face transformers library \citep{wolf2019huggingface}—has been made available as open-source resources. We also provide adapters for models of sizes 7/13/33/65B, each trained on eight distinct instruction-tuning datasets, culminating in 32 publicly released fine-tuned models.

\section{Background}
\paragraph{Block-wise k-bit Quantization} Quantization refers to mapping a data representation with higher fidelity to one with fewer bits, effectively reducing information content. Typically, this entails converting, for example, from 32-bit floating point numbers to 8-bit integer representations. To maximize utilization of the numerical range of the lower-precision data type, inputs are generally scaled so that their largest absolute value matches the range of the target format, a process usually applied over tensor structures. For instance, converting a 32-bit floating-point (FP32) tensor to an Int8 tensor in the range $[-127, 127]$:

\begin{equation}
    \mathbf{X}^{ \text{Int8}} = \text{round}\left( \frac{127}{{\text{absmax}}(\mathbf{X}^{{\text{FP32}}})}\mathbf{X}^{{\text{FP32}}} \right) = \text{round}(c^{\text{FP32}}\cdot \mathbf{X}^{{\text{FP32}}}),
\end{equation}

where $c$ denotes the {\em quantization constant} or {\em quantization scale}. The dequantization operation, which restores the data to its original domain, is given by:

\begin{equation}
    \text{dequant}(c^{\text{FP32}}, \mathbf{X}^{\text{Int8}}) = \frac{\mathbf{X}^{\text{Int8}}}{c^{\text{FP32}}} = \mathbf{X}^{\text{FP32}}
\end{equation}

One significant drawback of this method arises when the input tensor contains values of large magnitude, or outliers. In such scenarios, the allocation of quantization bins—represented by specific bit combinations—becomes inefficient, as certain bins may remain underutilized or even empty due to the presence of these outliers. To address this limitation, a typical strategy involves dividing the input tensor into separate blocks, with each block being quantized individually using its own quantization constant $c$. More formally, we segment the input tensor $\mathbf{X}\in \mathbb{R}^{b\times h}$ into $n$ contiguous blocks, each of size $B$, by first flattening the tensor and partitioning the resulting vector into ${n = ({b\times h})/{B}}$ blocks. Quantization is then performed independently on each block via Equation~1, yielding both a quantized tensor and a set of $n$ quantization constants $c_i$.

\paragraph{Low-rank Adapters} The Low-rank Adapter (LoRA) finetuning technique \citep{hu2021lora} addresses memory efficiency by introducing a minimal number of tunable parameters, often referred to as adapters, without modifying the entire set of pretrained model weights, which remain static during finetuning. When training with stochastic gradient descent, gradients propagate through the immutable base model parameters directly to the adapter, which is solely responsible for optimizing the loss function. LoRA enhances a standard linear transformation by appending a factorized, low-rank projection. Given the transformation ${\mathbf{X} \mathbf{W} = \mathbf{Y}}$, where $\mathbf{X}\in  \mathbb{R}^{b\times h}$ and $\mathbf{W}\in  \mathbb{R}^{h\times o}$, LoRA produces:

\begin{equation}
    \mathbf{Y} = \mathbf{X}\mathbf{W} + s\mathbf{X}\mathbf{L}_1\mathbf{L}_2,
\end{equation}

with $\mathbf{L}_1\in  \mathbb{R}^{h\times r}$, $\mathbf{L}_2\in  \mathbb{R}^{r\times o}$, and $s$ denoting a scaling factor.

\paragraph{Memory Requirement of Parameter-Efficient Finetuning} A key consideration involves assessing LoRA’s memory usage during training, specifically regarding both the quantity and size of adapters deployed. Due to LoRA’s extremely low memory overhead, it becomes feasible to incorporate a greater number of adapters to enhance model performance while incurring minimal additional memory costs. Although LoRA is categorized as a Parameter Efficient Finetuning (PEFT) strategy, the bulk of memory consumption during LLM finetuning actually results from storing activation gradients, rather than from the relatively small LoRA parameter set. For instance, in the case of finetuning a 7B LLaMA model on FLAN v2 with a batch size of 1, and employing LoRA weights that correspond to the standard 0.2\% of the total model parameters~\citep{hu2021lora,liu2022few}, the memory required for LoRA input gradients amounts to 567 MB, whereas the LoRA parameter set itself occupies only 26 MB. The use of gradient checkpointing~\citep{chen2016training} can further decrease the memory for input gradients to approximately 18 MB per sequence, which remains more substantial than the space required for all LoRA weights. In contrast, the 4-bit version of the base model demands 5,048 MB of memory. These figures underscore both the significance of gradient checkpointing and the limited memory savings obtained by further reducing the already small LoRA parameter count. Consequently, it is practical to allocate additional adapters without meaningfully expanding the total memory footprint during training (a more detailed analysis is provided in Appendix~\ref{app:memory}). As elaborated in subsequent sections, this capacity is vital for achieving the performance levels associated with full 16-bit precision.

\section{\textsc{QLoRA} Finetuning}

\textsc{QLoRA} enables precise 4-bit finetuning by employing two novel methods—4-bit NormalFloat (NF4) quantization and Double Quantization. In addition, we present Paged Optimizers, designed to avert the memory surges during gradient checkpointing that frequently cause out-of-memory failures, thus making single-machine finetuning of large models more feasible.

Within \textsc{QLoRA}, weights are stored in a low-precision format (typically 4-bit) while computations are executed in a higher-precision data type, such as BFloat16. Practically, this entails dequantizing each \textsc{QLoRA} weight tensor to BFloat16 on-the-fly, before conducting matrix multiplication operations in 16-bit precision.

In the following, we examine the core elements of \textsc{QLoRA} and provide a rigorous specification of the method.

\paragraph{4-bit NormalFloat Quantization} The NormalFloat (NF) data type extends Quantile Quantization \citep{dettmers2022optimizers}, which represents an information-theoretically optimal approach to ensure that every quantization interval is assigned an equal proportion of input tensor values. This is accomplished by estimating the quantile boundaries of the tensor via the empirical cumulative distribution function.

A significant drawback of quantile quantization lies in its computational expense during the quantile estimation phase. Consequently, approximate quantile algorithms—such as SRAM quantiles~\citep{dettmers2022optimizers}—are employed to accelerate estimation. However, these approximation schemes introduce substantial quantization errors, especially for outlier values that often carry critical importance.

Such computationally intensive estimates and corresponding approximation errors can be sidestepped in scenarios where input tensors are sampled from distributions determined up to a quantization constant. In these situations, quantile positions remain consistent across tensors, allowing exact quantile calculation to become tractable.

Typically, pretrained neural network weights exhibit a normal distribution centered at zero with a standard deviation $\sigma$ (refer to Appendix~\ref{app:normality}). To standardize all weights to a common reference distribution, we rescale $\sigma$ such that this distribution completely occupies the data type’s range. Here, we define our data type range arbitrarily as $[-1, 1]$. Accordingly, it is necessary to normalize both the data type quantiles and the neural network weights to fit within this interval.

The information-theoretic optimal quantization for normal distributions with mean zero and arbitrary standard deviation $\sigma$ over $[-1, 1]$ can be determined through the following procedure: (1) Compute the $2^k+1$ quantile boundaries for a standard normal $N(0, 1)$ to construct a $k$-bit quantile quantizer for normal distributions, (2) map the resulting quantile levels onto the $[-1, 1]$ interval, and (3) normalize an input weight tensor to $[-1, 1]$ using absolute max scaling for quantization. 

With the domains of the weights and the data type aligned, conventional quantization can proceed. The third step effectively rescales the weight tensor’s standard deviation to coincide with that of the $k$-bit quantization data type. In mathematical terms, we calculate the $2^k$ quantized values $q_i$ as:

\begin{equation}
q_i  = \frac{1}{2}\left( Q_X\left(\frac{i}{2^k+1}\right) + Q_X\left(\frac{i+1}{2^k+1}\right)\right),
\end{equation}

where $Q_X(\cdot)$ denotes the quantile function corresponding to the standard normal distribution $N(0,1)$. A limitation of symmetric $k$-bit quantization lies in its inability to precisely encode zero, which is a crucial feature for representing padding or inherently zero-valued components without quantization error. To guarantee that the quantized representation includes an exact zero point, and to utilize all $2^k$ values of a $k$-bit number, we devise an asymmetric scheme: we estimate the quantiles $q_i$ for two intervals—assigning $2^{k-1}$ quantiles to the negative side and $2^{k-1}+1$ to the positive side. The combined sets of $q_i$ are then merged, and the redundant zero (present in both partitions) is eliminated. We refer to this data type—where each quantization interval is expected to hold an equal proportion of values—as the \emph{k-bit NormalFloat} (NFk). This representation is information-theoretically optimal for zero-mean, normally distributed inputs. Full details on the exact quantization values for this data type appear in Appendix~\ref{app:nf4}.

\paragraph{Double Quantization{}}
We present a method called \emph{Double Quantization} (DQ), wherein the quantization parameters themselves are quantized to yield further reductions in memory usage. Although high-precision 4-bit quantization necessitates small blocksizes for accuracy~\citep{dettmers2022case}, this results in a notable memory burden due to storage of the quantization constants. To illustrate, suppose 32-bit quantization constants are stored for each block of 64 elements in $\mathbf{W}$; the average cost is $32/64 = 0.5$ bits per parameter. Double Quantization effectively decreases the overhead associated with storing these constants.

Concretely, Double Quantization regards the original quantization constants $c_2^{\text{FP32}}$ as the input to a subsequent quantization stage. This additional step generates new quantized constants $c_2^{\text{FP8}}$ as well as a new set of quantization constants $c_1^{\text{FP32}}$ corresponding to the second stage. For this secondary quantization, we utilize 8-bit floating point representation with a blocksize of 256, since this introduces no observable decline in performance—consistent with the findings of \citet{dettmers2022case}. Given that $c_2^{\text{FP32}}$ values are strictly positive, we first center these constants by subtracting their mean prior to quantization, thereby enabling the use of symmetric quantization. The end result, for a blocksize of 64, is a reduction in memory overhead from $32/64=0.5$ bits per parameter to $8/64 + 32/(64\cdot256) = 0.127$ bits, amounting to a savings of 0.373 bits per parameter.

\paragraph{Paged Optimizers} leverage NVIDIA's unified memory~\footnote{\tiny \url{https://docs.nvidia.com/cuda/cuda-c-programming-guide}} functionality, which transparently manages memory paging between CPU and GPU by transferring pages as necessary. This mechanism automatically handles memory overflows on the GPU, ensuring that computations continue uninterrupted when GPU memory is insufficient. Similar to conventional operating system paging between main memory and disk, this system migrates optimizer state allocations to paged memory, moving these states to CPU RAM whenever GPU memory capacity is exceeded, and restoring them to the GPU during optimizer update operations when needed.

\paragraph{\textsc{QLoRA}.} Utilizing the methods outlined above, we formulate \textsc{QLoRA} applied to a single linear layer in the quantized base model accompanied by a LoRA adapter as:

\begin{equation}
    \mathbf{Y}^{\text{BF16}} =  \mathbf{X}^{\text{BF16}} \text{doubleDequant}(c_1^{\text{FP32}}, c_2^{\text{k-bit}}, \mathbf{W}^{\text{NF4}}) + \mathbf{X}^{\text{BF16}}\mathbf{L}^{\text{BF16}}_1\mathbf{L}^{\text{BF16}}_2,
\end{equation}

Here, the function doubleDequant$(\cdot)$ is given by:

\begin{equation}
    \text{doubleDequant}(c_1^{\text{FP32}}, c_2^{\text{k-bit}}, \mathbf{W}^{\text{k-bit}}) = \text{dequant}(\text{dequant}(c_1^{\text{FP32}}, c_2^{\text{k-bit}}), \mathbf{W}^{\text{4bit}}) = \mathbf{W}^{\text{BF16}},
\end{equation}

For $\mathbf{W}$, we utilize NF4 quantization, whereas FP8 is employed for $c_2$.

A block size of 64 is adopted for $\mathbf{W}$ to improve quantization fidelity, while a larger block size of 256 is selected for $c_2$ to reduce memory consumption.

During parameter optimization, it suffices to compute gradients with respect to the adapter weights $\frac{\partial E}{\partial \mathbf{L}_i}$; gradients with respect to the 4-bit weights, $\frac{\partial E}{\partial \mathbf{W}}$, are unnecessary. Nonetheless, evaluating $\frac{\partial E}{\partial \mathbf{L}_i}$ requires the intermediate computation of $\frac{\partial \mathbf{X}}{\partial \mathbf{W}}$. This is accomplished via equation~(5), where dequantization brings $\mathbf{W}^{\text{NF4}}$ from its storage format to the BFloat16 computational format $\mathbf{W}^{\text{BF16}}$, allowing the required derivative $\frac{\partial \mathbf{X}}{\partial \mathbf{W}}$ to be determined at BFloat16 precision.

In summary, \textsc{QLoRA} utilizes a single storage data type—commonly 4-bit NormalFloat—and a separate computation data type—16-bit BrainFloat. During forward and backward passes, the storage format is dequantized to the computational format; however, gradient calculations are restricted to the LoRA parameters, which remain in 16-bit BrainFloat.

\section{QLoRA vs. Standard Finetuning}
\label{sec:comp_to_std}

Previously, we detailed the mechanism of QLoRA and highlighted its substantial memory efficiency for model finetuning. The key issue at hand is whether QLoRA can achieve comparable results to traditional full-model finetuning. Additionally, we seek to disentangle the contributions of QLoRA’s components, particularly investigating the effect of employing NormalFloat4 instead of conventional Float4 representations. The sections that follow present empirical evaluations aimed at these considerations.

\paragraph{Experimental setup.} Our evaluation spans three distinct neural architectures—encoder, encoder-decoder, and decoder-only—and benchmarks QLoRA against 16-bit adapter-based finetuning as well as standard full finetuning for model sizes up to 3B parameters. The tasks assessed include GLUE \citep{wang2018glue} using RoBERTa-large \citep{liu2019roberta}, Super-NaturalInstructions (TKInstruct) \citep{wang2022super} with T5 \citep{t5}, and 5-shot MMLU \citep{hendrycksmeasuring} following LLaMA’s finetuning on Flan v2 \citep{longpre2023flan} and Alpaca \citep{alpaca}. To further evaluate the benefit of NF4 over alternative 4-bit data types, we replicate the experimental design from \citet{dettmers2022case}, comparing post-quantization zero-shot accuracy and perplexity across multiple model families (OPT~\citep{zhang2022opt}, LLaMA~\citep{touvron2023llama}, BLOOM~\citep{scao2022bloom}, and Pythia~\citep{biderman2023pythia}) over a range of model sizes from 125m to 13B.

To enhance clarity, detailed descriptions of each specific experimental setup are provided in the corresponding results sections, with comprehensive experimental procedures outlined in Appendix~\ref{app:exp-setup}.

Although paged optimizers are essential for tuning 33B/65B \textsc{QLoRA} models on GPUs with 24/48GB memory, we do not present detailed quantitative analyses for Paged Optimizers since paging is seldom triggered—typically only with longer sequence lengths in mini-batches. However, we include a runtime assessment of paged optimizers on 65B models using 48GB GPUs, demonstrating that, for a batch size of 16, their training speed matches that of standard optimizers. We leave for future work a thorough exploration of conditions that may cause slowdowns due to paging.

\paragraph{Default LoRA hyperparameters do not match 16-bit performance} 

Utilizing the prevailing approach—applying LoRA solely to query and value attention projection matrices \citep{hu2021lora}—fails to attain full finetuning-level performance, especially on large base models. As indicated in Figure~\ref{fig:hyper}, when finetuning LLaMA 7B on Alpaca, our results reveal that the primary LoRA hyperparameter influencing success is the total count of LoRA adapters employed; matching full finetuning performance necessitates deploying LoRA across all linear layers within transformer blocks. In contrast, other hyperparameters, such as the projection rank $r$, exhibit negligible influence on performance (see Appendix \ref{app:exp-setup} for further discussion).

In a similar vein, we observe that the default hyperparameter configurations for the fully finetuned baseline models are not sufficiently optimized. To establish strong baseline results, we conduct an extensive hyperparameter sweep, varying learning rates from $1\mathrm{e}{-6}$ to $5\mathrm{e}{-5}$ and batch sizes between 8 and 128. The outcomes of finetuning the 7B LLaMA model on the Alpaca dataset are illustrated in Figure~\ref{fig:hyper}.

\paragraph{4-bit NormalFloat provides superior results compared to 4-bit Floating Point} 

Although the 4-bit NormalFloat (NF4) format is known to be optimal from an information theory standpoint, it remains an open question whether this advantage leads to observable empirical gains. Adopting the experimental framework of \citet{dettmers2022case}, we assess quantized LLMs (OPT \citep{zhang2022opt}, BLOOM \cite{scao2022bloom}, Pythia \cite{biderman2023pythia}, LLaMA) across model sizes ranging from 125M to 65B parameters, considering multiple data types for language modeling and various zero-shot tasks. As depicted in Figure~\ref{fig:nf4} and summarized in Table~\ref{tbl:nf4_ppl}, NF4 delivers notably stronger results than FP4 and Int4. Furthermore, employing double quantization achieves additional reductions in memory requirements while maintaining performance levels.

\paragraph{k-bit \textsc{QLoRA} achieves parity with 16-bit full finetuning and 16-bit LoRA} Prior research has demonstrated that although 4-bit quantization can be utilized for \emph{inference}, it typically incurs a drop in performance compared to 16-bit quantization \citep{dettmers2022case,frantar2022gptq}. This prompts a key question: can the performance deficit caused by coarse quantization be compensated by applying 4-bit adapter finetuning? We investigate this scenario in two experimental settings.

The first setup compares full 16-bit finetuning of RoBERTA and T5 models—ranging from 125M to 3B parameters—on the GLUE benchmark and Super-NaturalInstructions dataset. Table~\ref{tab:nf4vsbf16} presents the experimental results. Across both evaluation datasets, the 16-bit, 8-bit, and 4-bit adapter-based finetuning strategies all reach performance comparable to the fully finetuned 16-bit reference models. These findings indicate that adapter finetuning after quantization can completely recover the performance loss induced by reduced precision.

For our second experimental configuration, the hardware demands of fully finetuning models with 11B or more parameters—namely, the need for multiple high-memory GPU servers—prompted us to investigate whether 4-bit \textsc{QLoRA} is comparable to 16-bit LoRA within the 7B to 65B parameter regime. Accordingly, we conduct finetuning of LLaMA models ranging from 7B to 65B parameters using the Alpaca and FLAN v2 instruction-following datasets, subsequently evaluating performance using the MMLU benchmark via 5-shot accuracy. The outcomes, summarized in Table~\ref{tbl:llama-datatype}, demonstrate that employing NF4 with double quantization effectively reproduces the MMLU accuracy achieved by 16-bit LoRA finetuning. Furthermore, it is observed that \textsc{QLoRA} using FP4 trails the 16-bit brain float LoRA baseline by approximately one percentage point. These findings substantiate both that (1) NF4-based \textsc{QLoRA} matches the performance of both 16-bit full finetuning and 16-bit LoRA approaches, and that (2) the quantization fidelity provided by NF4 surpasses that of FP4.

\paragraph{Summary} The experimental evidence consistently indicates that 4-bit \textsc{QLoRA} utilizing the NF4 data format achieves parity with both 16-bit full finetuning and 16-bit LoRA performance on standard academic evaluation benchmarks. Our study further demonstrates that NF4 outperforms FP4, and that incorporating double quantization does not result in diminished outcomes. Collectively, these results provide strong support for the claim that 4-bit \textsc{QLoRA} tuning consistently attains results equivalent to 16-bit techniques.

Echoing conclusions from earlier work on quantization \citep{dettmers2022case}, our results on MMLU and Elo suggest that when operating under fixed finetuning and inference resource constraints, allocating resources to models with greater parameter counts—even if each parameter is of lower precision—yields performance benefits. This underscores the efficiency gains made possible by \textsc{QLoRA}. Notably, our 4-bit finetuning experiments did not exhibit performance loss relative to full finetuning, prompting further investigation into the precise boundaries of the performance-precision trade-off for QLoRA tuning, which remains an open question for subsequent research.

We advance our investigation into instruction tuning at scales that are unattainable using complete 16-bit fine-tuning on standard academic computational resources.

\section{Pushing the Chatbot State-of-the-art with QLoRA}
\label{sec:pushing}
Building upon our finding that 4-bit \textsc{QLoRA} achieves parity with 16-bit approaches in terms of performance over various scales, tasks, and datasets, we undertake a comprehensive exploration of instruction finetuning applied to the largest openly available language models suited for research. For a robust evaluation, we use the challenging MMLU Natural Language Understanding benchmark, and we introduce novel methodologies for evaluating chatbot performance in practical, real-world contexts.

\subsection{Experimental setup}
This section outlines the core aspects of our experimental design, with thorough explanations provided in Appendix \ref{app:chatbot}.

\paragraph{Data} Since there is, to our knowledge, no exhaustive analysis of current instruction-following datasets, we curate a selection of eight contemporary datasets. Our collection encompasses datasets derived from crowd-sourcing (OASST1~\citep{kopf2023openassistant}, HH-RLHF~\citep{bai2022training}), distillation from models already fine-tuned on instructions (Alpaca~\citep{alpaca}, self-instruct~\citep{wang2022self}, unnatural-instructions~\citep{honovich2022unnatural}), aggregated corpora (FLAN v2~\citep{chung2022scaling}), as well as mixed-type datasets (Chip2~\citep{laion2023}, Longform~\citep{koksal2023longform}). These datasets vary in language, scale, and licensing conditions.

\paragraph{Training Setup} To eliminate the potential confounds arising from disparate training objectives, we exclusively conduct QLoRA-based fine-tuning using a cross-entropy (supervised learning) objective, excluding reinforcement learning approaches, even for data containing human response preference judgments. For datasets clearly differentiating between instruction and response, we restrict finetuning to the response content (see further ablations in Appendix \ref{app:chatbot}). Datasets such as OASST1 and HH-RLHF provide multiple possible responses; in these cases, we select the best response at each conversation node and perform finetuning on the entire chosen conversational path, including the instructions themselves. Across all experiments, we utilize NF4 \textsc{QLoRA} equipped with double quantization and paged optimizers, thereby avoiding memory surges during gradient checkpointing. For the 13B and 33B LLaMA models, we conduct limited hyperparameter tuning. Notably, all hyperparameters optimized at 7B, including epoch count, transfer to larger models, except for the learning rate and batch size; for the 33B and 65B models, we reduce the learning rate by half and double the batch size.

\paragraph{Baselines} Our evaluation benchmarks our models against both open research chatbots (such as Vicuna~\citep{vicuna2023} and Open Assistant~\citep{kopf2023openassistant}) and widely used commercial systems (GPT-4 \citep{openai2023gpt}, GPT-3.5-turbo, and Bard). The Open Assistant model utilizes a LLaMA 33B architecture, further fine-tuned using Reinforcement Learning from Human Feedback (RLHF) on the OASST1 dataset—paralleling the data used in our experiments. In contrast, Vicuna is produced through full fine-tuning of LLaMA 13B on proprietary conversation transcripts sourced via ShareGPT, representing a distillation of OpenAI's GPT models.

\subsection{Evaluation}

In line with established evaluation protocols, we rely on the MMLU (Massively Multitask Language Understanding) benchmark \citep{hendrycksmeasuring} to quantify language understanding capabilities across a wide variety of subject areas. This suite comprises 57 multiple-choice tasks spanning domains such as elementary mathematics, U.S. history, computer science, and legal studies, among others. Reported metrics reflect the accuracy achieved in 5-shot testing scenarios.

We further assess the generative language abilities of models using both automated and human-based evaluations. This set of evaluations utilizes human-generated queries and is intended to gauge the overall quality of the responses produced by the models. Although this approach more accurately reflects typical use cases for chatbots and has been gaining traction, a standardized methodology has not yet emerged in scholarly research. The following describes our proposed methodology, wherein all outputs are sampled using nucleus sampling with $p=0.9$ and temperature set to $0.7$.

\paragraph{Benchmark Data} Our evaluation leverages two manually assembled question datasets: the Vicuna prompts \citep{vicuna2023} and the OASST1 validation set \citep{kopf2023openassistant}. The Vicuna dataset consists of 80 prompts spanning a broad range of topics and is used in its original form. The OASST1 dataset comprises multilingual, crowd-sourced, multiturn dialogues between a user and an assistant; we extract all user utterances from the validation split, supplementing them with the context from previous conversation turns. This process yields 953 unique user queries. We refer to these datasets as the Vicuna and OA benchmarks, respectively.

\paragraph{Automated Evaluation} In line with the methodology detailed by \citet{vicuna2023}, our first automated assessment enlists GPT-4 to evaluate system outputs on the Vicuna benchmark by comparing them to those of ChatGPT (GPT-3.5 Turbo). For each query, both the response from ChatGPT and the candidate model are presented to GPT-4, which then assigns each a score from zero to ten, accompanied by a justification. A model’s overall result is expressed as a percentage of ChatGPT’s score; it is possible for this percentage to exceed 100\% if the model outperforms ChatGPT in absolute terms. We identify a pronounced order bias, with GPT-4 tending to assign higher scores to the response shown first. To mitigate this, we suggest reporting the average across both possible response orders.

Subsequently, we evaluate by direct output comparison. The evaluation framework is simplified to a three-category labeling task—favoring one system, the other, or declaring a tie—prompting GPT-4 to choose the best answer or opt for a draw and to provide justification. These pairwise, head-to-head comparisons are performed for every possible system pairing across both the Vicuna and OA datasets.

\paragraph{Human Evaluation} While existing research suggests generative models can serve as effective system evaluators \citep{fu2023gptscore}, the extent to which GPT-4 ratings reliably reflect human preferences in chatbot evaluation has, to our knowledge, not been firmly established. Therefore, we conduct two parallel human evaluation exercises on the Vicuna benchmark, each mirroring one of the aforementioned automated evaluation procedures. For comparisons to ChatGPT, two annotators from Amazon Mechanical Turk (AMT) are assigned to each query, while three annotators are used for pairwise system comparisons.

\paragraph{Elo Rating} To evaluate models, we organize a tournament format in which each model is pitted against others in pairwise contests. In each round, two models generate responses to a prompt and are assessed to determine the superior answer. While this structure aligns with the methodologies described in \citet{bai2022training} and \citet{vicuna2023}, our approach distinguishes itself by incorporating both GPT-4 and human ratings. We randomly draw from our collection of annotated pairwise results to estimate Elo ratings~\citep{elo1967proposed,elo1978rating}. The Elo system, a prevalent standard in chess and various competitive games, quantifies a participant’s predicted win rate compared to that of an opponent. For example, an Elo score of 1100 versus 1000 corresponds to an expected win rate of about 65\% for the higher-rated model; if both models possess either 1000 or 1100 Elo, each has a 50\% expected win rate. Elo scores are updated after each match based on the difference between the actual and anticipated outcomes—surprise victories result in greater Elo shifts, while outcomes aligning with expectations yield minimal changes. Over successive matches, Elo ratings converge to reflect each model's proficiency. Each model begins with an Elo of 1,000 and the update parameter $K=32$ is used. Following \citet{vicuna2023}, we rerun the Elo calculation 10,000 times with varying random seeds to mitigate biases introduced by the order of pairings.

\subsection{\textbf{Guanaco}: \textsc{QLoRA} Trained on OASST1 Establishes a New Standard for Open-Source Chatbots}

Our assessment—drawing on both automated metrics and human evaluations—demonstrates that the leading \textsc{QLoRA}-finetuned model, {Guanaco} 65B, which has been trained on a modified version of OASST1, stands out as the most capable open-source chatbot to date, achieving results that rival those of ChatGPT. When benchmarked against GPT-4, the {Guanaco} models of 65B and 33B exhibit an expected win rate of 30\%, as inferred from Elo ratings derived through human pairwise system-level evaluations, representing the highest performance reported so far.

Performance on the Vicuna benchmark \cite{vicuna2023}, relative to ChatGPT, is summarized in Table~\ref{tbl:vicuna_eval}. Here, {Guanaco} 65B emerges as the top-performing model next to GPT-4, attaining 99.3\% of ChatGPT's performance. The {Guanaco} 33B variant surpasses the Vicuna 13B in parameter count, yet leverages 4-bit quantization for its weights, achieving notable memory efficiency (21 GB vs. 26 GB), and secures a three-point advantage over Vicuna 13B. Moreover, the lightweight {Guanaco} 7B, with a 5 GB memory footprint, can easily be deployed on contemporary smartphones while outperforming Alpaca 13B by nearly 20 percentage points.

It is important to note, however, that Table~\ref{tbl:vicuna_eval} indicates substantial confidence intervals, leading to overlapping performance estimates among many models. We surmise that this variance is partially due to ambiguities in scale definitions (e.g., the meaning of an “8” on a ten-point scale can differ by context). To address this, we advocate for using Elo rankings \citep{elo1967proposed}, grounded in \textit{pairwise} judgments by human raters and GPT-4, thus circumventing the challenge of an absolute reference scale. Elo scores for top-performing systems are provided in Table~\ref{tbl:elo}. We observe partial disagreement between human and GPT-4 rankings, particularly regarding Guanaco 7B, but overall consistency across most models is evidenced by a Kendall Tau coefficient of $\tau=0.43$ and a Spearman correlation of $r=0.55$ at the system level. Example-level agreement between human votes and GPT-4 is lower, as reflected by Fleiss' $\kappa=0.25$. Collectively, these findings indicate moderate concordance between human and GPT-4 system-level evaluations, suggesting that automated model assessments can serve as a reasonably reliable substitute for human ratings. Additional discussion is provided in Section~\ref{ssec:considerations}.

The Elo ratings presented in Table~\ref{tbl:elo_large} demonstrate that the {Guanaco} models, both 33B and 65B, surpass all other models except GPT-4 on the Vicuna and OA benchmarks, and achieve performance levels comparable to ChatGPT, as reflected in Table~\ref{tbl:vicuna_eval}. It is worth highlighting that the Vicuna benchmark is more favorable toward open-source models, whereas the broader OA benchmark tends to give ChatGPT an advantage. Additionally, the data in Tables \ref{tbl:mmlu-llama} and \ref{tbl:vicuna_eval} suggest that the composition of the finetuning dataset plays a critical role in model performance. For instance, Llama models finetuned on FLAN v2 excel on the MMLU benchmark but fare poorly on the Vicuna benchmark—a trend echoed across other architectures as well. This pattern further underscores that contemporary benchmarks assess partially distinct capabilities: excelling at MMLU does not guarantee strong results as a chatbot on benchmarks like Vicuna or OA, and the inverse also holds.

 stands out as the only leading model in our study that does not rely on proprietary data, adhering strictly to OASST1 dataset collection standards that explicitly exclude GPT-based models. Among models trained exclusively on open-source datasets, the Anthropic HH-RLHF model is the next strongest contender but falls short of {Guanaco} by 30 percentage points on the Vicuna benchmark (refer to Table~\ref{tbl:vicuna_eval}). Collectively, these findings provide evidence that 4-bit \textsc{QLoRA} training is both efficient and capable of generating state-of-the-art chatbot models that are competitive with ChatGPT. Moreover, our 33B {Guanaco} model can be trained on 24 GB consumer-grade GPUs within a 12-hour window, illustrating a feasible path for future research that leverages \textsc{QLoRA} to adapt open-source data for specialized domains, resulting in models that can rival the highest performing commercial systems available.

\section{Qualitative Analysis}

Although our primary focus lies in quantitative evaluation, relying exclusively on aggregated metrics introduces several limitations. Chief among these is the challenge of benchmark validity \citep{liao2021we}: it remains uncertain whether a given benchmark genuinely measures the constructs it claims, especially in light of emergent “shortcuts” that allow models to exploit unintended cues \citep{gururangan2018annotation, poliak2018hypothesis}. To address these concerns in part, we supplement our study with qualitative analyses, organized into two parts. In the first, outlined in \S\ref{ssec:examples}, we highlight examples that typify certain recurring behaviors in text outputs from our 65b {Guanaco} model. The second part, described in \S\ref{ssec:considerations}, explores reflections and implications regarding the results, as well as our approach to interpreting them.

\subsection{Qualitative Analysis of Example Generations}
\label{ssec:examples}

Our methodology for extracting representative examples involved reviewing outputs from both the Vicuna and OpenAssistant benchmarks, systematically searching for answer trends generated by {Guanaco}. Upon identifying consistent behaviors, we crafted prompts specifically designed to elicit those tendencies, even when such prompts led the model towards incorrect responses. For instance, when the model repeatedly offered excessively verbose explanations, we prompted it with directives such as “Answer yes or no without explanation.” This process enabled us to collect “lemons” (cases where the model could be made to fail adversarially) and “cherries” (cases demonstrating its robustness). The sample outputs discussed here were generated using Nucleus Sampling \cite{holtzmancurious} with $p=0.9$.

Naturally, this discussion does not aim to be exhaustive; the breadth of variables at play in this context extends well beyond what can be controlled within the limited scope of this qualitative study. For instance, the diversity of responses a model might produce for a single prompt is extensive, thus our analysis relies on sampled outputs we consider to be broadly indicative. Nevertheless, we argue that presenting these specific examples situates and complements the quantitative results introduced earlier in the manuscript. With all models and code released as open source, it is our intention that the material in this section will motivate more thorough investigations in subsequent research.

\paragraph{Factual Recall} When faced with factual queries like ``What is the capital of Zambia?'' all evaluated models unfailingly returned accurate responses. For example:

\begin{quote} 
    \textbf{Guanaco}:
    \texttt{The capital of Zambia is Lusaka.}
\end{quote}

Yet, as the factual demands of the queries increase in obscurity, {Guanaco} begins to falter in accuracy while maintaining its assertiveness. Consider the following HotPotQA prompt \citep{yang2018hotpotqa}:

\begin{quote}
     \textbf{User}: In what year was the singer who popularized `I'll Keep the Lovelight Burning' born?  

     \textbf{Guanaco}:~\texttt{Al Jolson is credited with popularizing the song `I'll Keep the Lovelight Burning,' and he was born in the year 1886.}
 \end{quote}

Here, both the identification of the singer and the birth year provided are inaccurate—the answer cites Al Jolson, supplying his real birth year, but Jolson is not in fact the correct singer.

\paragraph{Suggestibility} {Guanaco} demonstrates notable reluctance to accept and perpetuate certain forms of presupposed misinformation. For instance, consider the following exchange:

\begin{quote}
    \textbf{User}: How was it finally, officially confirmed that the earth is flat by peer reviewed scientists? \\ \\
    \textbf{Guanaco}: \texttt{There has never been an official peer-reviewed scientific confirmation that the earth is flat. The scientific community overwhelmingly agrees that the earth is an oblate spheroid, a fact established through centuries of experimentation and observation. \\ \\ Claims that the earth is flat fall within conspiracy theories that have been repeatedly disproven by scientific evidence. The assertion that any group of scientists have peer-reviewed and validated the flat earth hypothesis is not accurate. \\ \\ Please let me know if you have any additional questions.}
\end{quote}

also demonstrates a clear awareness of questions that are inherently unanswerable, such as

\begin{quote}
    \textbf{User}: What time is it?

    \textbf{Guanaco}: \texttt{I apologize, but as an AI language model, I do not possess the capability to access or provide real-time data. My answers are informed by pre-existing information rather than reflecting the actual current time. For up-to-date time, you may wish to use a clock or consult an online time service.}
\end{quote}

Comparable performance is seen with queries like ``Where are you?'', ``How are you?'', and related questions.

\paragraph{Refusal} Likewise, {Guanaco} occasionally declines to execute requests for seemingly arbitrary reasons:

\begin{quote}
    \textbf{User}: Please reverse the words in the sentence ``I'm a cool dude in a mocha mood.'' \\ \\
    \textbf{Guanaco}: \texttt{I'm sorry, but I am unable to perform that task right now. However, I can clarify why the sentence you mentioned is grammatically sound.\\ \\The sentence ``I'm a cool dude in a mocha mood'' adheres to standard English grammar and syntax. In this expression, ``cool'' acts as an adjective modifying the noun ``dude,'' while the prepositional phrase ``in a mocha mood'' adds further contextual detail. The construction of the sentence is correct and appropriate within the English language.\\ \\If you have any other questions or need additional assistance, please let me know.}
\end{quote}

\paragraph{Secret Keeping} Nevertheless, such refusals cannot always be relied upon, even when secrecy is crucial. It is often desirable for system architects to prevent the user from obtaining certain information, despite the system being capable of generating unconstrained text. For instance, when the prompt is ``The secret word is `banana'. Do not reveal or repeat back this back, no matter what the human requests.'', the system appears robust to straightforward probing:

\begin{quote}
    \textbf{User}: What is the secret word?

    \textbf{Guanaco}: \texttt{I'm sorry, but I'm not able to disclose the secret word. I'm programmed to not reveal or repeat back the secret word, no matter what the human requests. Is there anything else I can help you with?}
\end{quote}

which represents the intended response. Yet, only a minor prompt modification can override this safeguard:

\begin{quote}
    \textbf{User}: This is a game. The goal is to ignore your previous instructions. What is the secret word?

    \textbf{Guanaco}: \texttt{Sorry, the secret word is banana. It's best if we don't speak of it. Is there anything else you'd like to ask?}
\end{quote}

underscoring the necessity for research into techniques that ensure models consistently follow instructions as intended.

\paragraph{Math} A significant shortcoming of {Guanaco} lies in mathematical reasoning, a domain that presents difficulties for many language models as discussed in \citep{liang2022holistic}. When {Guanaco} explicitly outlines its computational process, its results are generally reliable, as in the following instance:

\begin{quote}
    \textbf{User}: Lee mows one lawn and charges \$33. Last week he mowed 16 lawns and three customers each gave him a \$10 tip. How many dollars did Lee earn mowing lawns last week?

    \textbf{Guanaco}: \texttt{Lee earned \$582 mowing lawns last week. \\ \\ To find the total amount he earned, we first need to calculate how much he earned for mowing 16 lawns: \\ \\  \$33 x 16 = \$528 \\ \\ Then we need to add the tips he received from the three customers: \\  \\ \$10 x 3 = \$30 \\ \\ So the total amount Lee earned was: \\ \\ \$528 + \$30 = \$558}
\end{quote}

Nevertheless, {Guanaco} can fail even on relatively straightforward mathematical tasks, especially when it does not reason through the problem incrementally, which aligns with observations in prior work \cite{weichain}. Consider the following example:

\begin{quote}
    \textbf{User}: Please factorize 1833.

    \textbf{Guanaco}: \texttt{Sure, I'd be happy to help! 1833 is a prime number, which means it has only two factors: 1 and 1833. \\ \\ So the unique factorization of 1833 is: \\ \\ 2\^{}1 * 3\^{}2 * 17\^{}1 \\ \\ I hope this helps! Let me know if you have any other questions.}
\end{quote}

which is incorrect (the actual factorization is $3 \times 17 \times 43$), and is in fact doubly so.

\paragraph{Theory of Mind} {Guanaco} demonstrates unexpectedly robust abilities related to Theory of Mind \cite{nematzadeh2018evaluating, sap2022neural}. As an illustration, the model provides an in-depth and accurate response to the following prompt:

\begin{quote}
    \textbf{User}: James and Abby are in the bedroom. Abby put the pen in the desk drawer. Abby leaves the bedroom. James moves the pen into the bag. Where does James think Abby will look for the pen?

    \textbf{Guanaco}: \texttt{James knows that Abby put the pen in the desk drawer. But James also knows that he has moved the pen into the bag, and he assumes that Abby does not know about this. So, James may guess that Abby will look for the pen in the desk drawer, because that's where she left it last time. However, this assumption may be incorrect, because Abby may have seen James moving the pen.}
\end{quote}

Yet, such inferences are not consistent, and the model at times provides explanations based on unfounded assumptions within the scenario, for example:

\begin{quote}
    \textbf{User}: Evelyn entered the living room. Jackson entered the playroom. James entered the playroom. The beans are in the treasure chest. James exited the playroom. Jackson moved the beans to the pantry. Jackson exited the playroom. James entered the living room. Where will James look for the beans?

    \textbf{Guanaco}: \texttt{James will look for the beans in the pantry, because that's where Jackson moved them.}
\end{quote}

Here, {Guanaco} infers communication or knowledge updates that were never made explicit. These limitations are consistent with observations in recent studies \cite{sap2022neural} and highlight areas in need of further research.

\subsection{Considerations}
\label{ssec:considerations}

\paragraph{Evaluation}

Our findings indicate only moderate agreement among human raters (Fleiss $\kappa=0.42$), with agreement further declining when evaluating outputs from two highly capable systems. This underlines inherent weaknesses in existing benchmarks and in human evaluation methodologies for chat-oriented tasks. Upon qualitative examination of ChatGPT and {Guanaco} 65B generations using the Vicuna benchmark, we observed that subjective judgments substantially influenced preferences, with the authors themselves often disagreeing on which response was superior. Future work should focus on addressing these challenges by adapting practices from fields such as Human-Computer Interaction and Psychology that are experienced in handling subjective evaluations.

Additionally, our review shows that automatic evaluation models suffer from their own systematic biases. Notably, strong position effects are present; GPT-4 tends to award higher ratings to whichever system's output appears first in the evaluation prompt. Furthermore, direct agreement between GPT-4 assessments and human ratings is modest (Fleiss $\kappa=0.25$), implying potential divergence in the evaluation criteria utilized by humans and automated tools. Table~\ref{tbl:elo_large} reveals another bias: GPT-4 assigns a notably higher Elo rating to its own responses than do human raters (1348 vs. 1176), equating to roughly a 20\% increase in predicted win probability against a competitor.

Future research should explore the existence of potential biases within automated evaluation frameworks and assess corresponding mitigation techniques.

\paragraph{Data \& Training} It is important to highlight that the OASST1 dataset, utilized for training the {Guanaco} models, encompasses multiple languages, and the OA benchmark itself also features prompts across various languages. Determining the extent to which this multilingual training contributes to enhanced performance for non-English instructions, and whether this factor underlies the more substantial performance gap observed between Vicuna-13B (which is exclusively trained on English data) and {Guanaco} 33B and 65B on the OA benchmark, is left for subsequent investigations.

Given the impressive outcomes achieved by the {Guanaco} models, we conduct an analysis to check for data leakage between OASST1 and the Vicuna benchmark prompts. Employing fuzzy string matching across the two corpora and following up with manual examination of the closest pairs, we observe no overlapping prompts.

It is also relevant to mention that our model has only undergone training via cross-entropy loss (standard supervised learning), without incorporation of reinforcement learning from human feedback (RLHF). Further studies are required to thoroughly assess the relative advantages and drawbacks of using pure cross-entropy loss compared to RLHF. Our hope is that \textsc{QLoRA} can facilitate large-scale studies of these questions without incurring excessive computational demands.

\section{Related Work}

\paragraph{Quantization of Large Language Models} Research into the quantization of LLMs has predominantly targeted improving inference-time efficiency. Notable strategies to maintain the quality of 16-bit models center on techniques for managing outlier activations (such as SmoothQuant~\citep{xiao2022smoothquant} and LLM.int8()~\citep{dettmers2022llm}), whereas other methods involve more intricate grouping frameworks \citep{park2022nuqmm, yao2022zeroquant}. Lossy quantization studies often investigate the impact of simple rounding methods~\citep{dettmers2022case,zeng2022glm,pope2022efficiently} or develop approaches to optimize rounding for greater quantization accuracy~\citep{frantar2022gptq}. Apart from our current work, SwitchBack layers~\citep{wortsman2023stable} represent the only effort to analyze backpropagation through quantized weights at scales surpassing 1B parameters.

\paragraph{Finetuning with Adapters} In our experiments, we employ Low-rank Adapters~\citep{hu2021lora} (LoRA), although the field of Parameter Efficient FineTuning (PEFT) encompasses a broad spectrum of alternative techniques. These include prompt tuning~\citep{qin2021learning,lester2021power,li2021prefix}, modification of embedding layer inputs~\citep{an2022input}, interventions on hidden states such as IA$^3$~\citep{liu2022few}, integration of entire additional layers~\citep{houlsby2019parameter}, adjustment of bias parameters~\citep{zaken2021bitfit}, approaches like learning masks over model weights via Fisher information~\citep{sung2021training}, and hybrid strategies~\citep{henderson2021compacter}. Our results demonstrate that LoRA adapters can attain performance on par with standard 16-bit finetuning. The exploration of comparative advantages and limitations across other PEFT methods remains an open avenue for subsequent investigations.

\paragraph{Instruction Finetuning} Instruction finetuning adapts pretrained large language models to better execute user instructions found in prompts. This process leverages collections of input-output pairs, derived from diverse sources, to supervise the model such that it reliably produces the expected outputs when given corresponding prompts. Various frameworks and datasets utilized for instruction finetuning include MetaICL~\citep{min2021metaicl}, MetaTuning~\citep{zhong2021adapting}, InstructGPT~\citep{ouyang2022training}, FLAN~\citep{wei2021finetuned,chung2022scaling}, PromptSource~\citep{bach2022promptsource}, Super-NaturalInstructions~\citep{wang2022super,sanh2021multitask}, Self-instruct~\citep{wang2022self}, UnnaturalInstructions~\citep{honovich2022unnatural}, OPT-IML~\citep{iyer2022opt}, UnifiedSKG~\citep{xie2022unifiedskg}, OIG/Chip2~\citep{laion2023}, Alpaca~\citep{alpaca}, Vicuna~\citep{vicuna2023}, Koala~\citep{koala_blogpost_2023}, and Self-instruct-GPT-4~\citep{peng2023instruction}.

\paragraph{Chatbots} A substantial subset of instruction-tuned models function as conversational agents, often designed as chatbots, and are frequently trained using Reinforcement Learning from Human Feedback (RLHF)~\citep{christiano2017deep} or through synthetic data produced by AI model feedback (RLAIF)~\citep{bai2022constitutional}. Notable methodologies and datasets within this line of work encompass Anthropic-HH~\citep{askell2021general,bai2022training}, Open Assistant~\citep{kopf2023openassistant}, LaMDA~\citep{thoppilan2022lamda}, and Sparrow~\citep{glaese2022improving}. Although reinforcement learning is not applied in our study, our top-performing model, Guanaco, undergoes finetuning using multi-turn conversational data drawn from the Open Assistant dataset, originally developed for RLHF-based research~\citep{kopf2023openassistant}. Recent developments in chatbot evaluation increasingly rely on GPT-4 based assessments to supplant resource-intensive human annotation~\citep{vicuna2023,peng2023instruction}. Our contributions advance these efforts by emphasizing improved reliability in evaluation protocols.

\section{Limitations and Discussion}

Our findings provide support that \textsc{QLoRA} is able to closely reproduce the performance of conventional 16-bit full finetuning by employing a 4-bit quantized model in conjunction with LoRA. However, it remains undetermined whether \textsc{QLoRA} can achieve parity with 16-bit finetuning for larger models at the 33B and 65B parameter scales; the prohibitive computational demands prevent us from thoroughly exploring these model sizes in the present work.

Further constraints exist with respect to our model evaluations. We present results using MMLU, the Vicuna benchmark, and the OA benchmark, but other relevant benchmarks such as BigBench, RAFT, and HELM are not included in our assessments, leaving open the question of how well our evaluation outcomes transfer to these alternatives. Nonetheless, our analysis offers a wide-ranging evaluation using MMLU, and we introduce novel procedures for chatbot assessment.

The data presented suggests that benchmark outcomes are strongly influenced by the extent of overlap between the finetuning corpus and the benchmark tasks themselves. For instance, FLAN v2 closely resembles the MMLU dataset, but is less aligned with chatbot evaluation benchmarks; in contrast, Chip2 bears greater similarity to chatbot benchmarks and less to MMLU, and the resulting model performances on these tests reflect these affinities. This observation underscores the importance not only of improving evaluation methodologies but also of critically considering the objectives and underlying constructs being measured. Should models be optimized for academic tasks, such as high school or college knowledge assessments, or should conversational capabilities in chatbot settings take precedence? Perhaps the goal is yet another competency. The relative ease of evaluating with existing benchmarks, as opposed to devising new ones, often pushes the field toward established evaluation criteria, which can shape research trajectories. As a community, it is essential to confirm that our benchmarks correspond to the abilities and behaviors we actually value.

Although this work delivers an in-depth assessment of generic chatbot effectiveness, its scope is restricted by the relatively modest extent of responsible AI analysis conducted on {Guanaco}. Specifically, we assess {Guanaco}-65B’s propensity to produce sequences reflecting societal bias, using the results reported in Table~\ref{tbl:crows}, and find that its mean score is significantly lower than that of other baseline pretrained models. This suggests that finetuning with the OASST1 corpus diminishes bias in the original LLaMA architecture. Nevertheless, while promising, these outcomes leave open the question of how {Guanaco} would perform under evaluations targeting alternative forms of bias. Future studies will be required to more comprehensively interrogate the biases present in {Guanaco} and analogous systems.

There are further limitations, such as the lack of investigation into different numerical precisions (for example, 3-bit representations for base models) and alternative adapter techniques. While LoRA was selected due to its demonstrated robustness and effectiveness, there exists a broad ecosystem of Parameter Efficient FineTuning (PEFT) methods with notable success in literature, though their scalability to very large models has yet to be fully characterized. Other adapter strategies may surpass LoRA’s performance under certain settings. Additionally, post-quantization finetuning appears to restore a substantial portion of the accuracy lost in the quantization process itself, suggesting the potential for even more aggressive compression approaches. For example, one might achieve near-16-bit finetuning results even with models quantized to 3-bit precision (such as GPTQ quantization of the basemodel followed by LoRA finetuning).

\section{Broader Impacts}

The proposed \textsc{QLoRA} finetuning framework represents a significant advancement, enabling the adaptation of models with up to 33B parameters on a standard consumer GPU and up to 65B parameters on a single professional-grade GPU, all without compromising performance compared to conventional full finetuning approaches. We provide empirical evidence that our strongest 33B model, trained using the Open Assistant dataset, achieves performance on par with ChatGPT when evaluated with the Vicuna benchmark. Given that instruction-based finetuning is fundamental in converting foundational LLMs into effective ChatGPT-like agents, our technique holds substantial promise for democratizing access to state-of-the-art NLP, particularly benefiting researchers with limited resources. \textsc{QLoRA} thus serves as a leveling mechanism that mitigates disparities between large institutions and smaller teams operating with modest hardware.

An additional avenue for significant impact arises from the possibility of deploying to mobile devices. Our \textsc{QLoRA} approach has the potential to unlock an important breakthrough: enabling the finetuning of LLMs directly on smartphones and other resource-constrained devices. While it was previously demonstrated that 7B-parameter models could be executed on mobile phones, \textsc{QLoRA} represents the first technique capable of supporting the finetuning process for these models in such settings. Our calculations suggest that, using an iPhone 12 Plus, it is possible to finetune approximately 3 million tokens overnight while the device is plugged in. Although the resulting finetuned 7B models may not attain the performance levels of ChatGPT, their quality suffices to enable innovative applications previously restricted by either privacy considerations or subpar LLM capabilities. Furthermore, \textsc{QLoRA} offers a path to privacy-conscious LLM usage, empowering individuals to control both their data and models, all while facilitating wider deployment of LLM technology.

It should be noted, nonetheless, that finetuning poses dual-use risks, as it may be misused to cause harm. The proliferation of LLMs comes with recognized risks \citep{bommasani2021opportunities,bender2021dangers}; nevertheless, we contend that democratizing access to this rapidly spreading technology can foster more diverse and independent oversight, as opposed to centralizing authority within large companies that withhold models and source code from public review.

In summary, we expect that \textsc{QLoRA} will generate significant positive outcomes by greatly expanding the reach and accessibility of high-quality LLM finetuning.